% Part: applied-modal-logics
% Chapter: epistemic-logic
% Section: language-epistemic-logic

\documentclass[../../../include/open-logic-section]{subfiles}

\begin{document}

\olfileid{aml}{el}{pal}

\olsection{सार्वजनिक घोषणा तर्कशास्त्र}

गतिशील ज्ञानविषयक तर्कशास्त्रांमुळे कर्त्यांचे ज्ञान
काळानुसार किंवा नवी माहिती मिळाल्यावर कसे बदलते हे
दर्शवता येते. यांपैकी अनेक तर्कशास्त्रे ज्ञानातील बदल
माहितीपर \emph{घटना} किंवा \emph{अद्यतने} यांद्वारे
दर्शवतात. अद्यतनाचा सर्वांत मूलभूत प्रकार म्हणजे
सार्वजनिक घोषणा: एखादे सूत्र सत्य असल्याची घोषणा
होते आणि सर्व कर्ते ती एकत्रितपणे पाहतात. यासाठी
भाषेचा पुढीलप्रमाणे विस्तार करतो.

\begin{defn}
$G$ हा कर्तृचिन्हांचा संच असू द्या. सार्वजनिक घोषणा
असलेल्या अनेक कर्त्यांच्या ज्ञानविषयक तर्कशास्त्राच्या
मूलभूत भाषेत पुढील गोष्टी असतात:
\begin{enumerate}
  \tagitem{prvFalse}{!!{falsity} दर्शवणारा विधानीय स्थिरांक~$\lfalse$.}{}
  \tagitem{prvTrue}{!!{truth} दर्शवणारा विधानीय स्थिरांक~$\ltrue$.}{}
  \item !!{propositional variable}s चा !!{denumerable}s संच: $\Obj
    p_0$, $\Obj p_1$, $\Obj p_2$, \dots
  \item विधानीय संयोजके: \startycommalist
  \iftag{prvNot}{\ycomma $\lnot$ (नकार)}{}%
  \iftag{prvAnd}{\ycomma $\land$ (संयोग)}{}%
  \iftag{prvOr}{\ycomma $\lor$ (विकल्पयोग)}{}%
  \iftag{prvIf}{\ycomma $\lif$ (!!{conditional})}{}%
  \iftag{prvIff}{\ycomma $\liff$ (!!{biconditional})}{}.
  \item $a \in G$ असताना ज्ञान-संकारक $\Knows_a$.
  \item $!B$ हे !!a{formula} असताना सार्वजनिक घोषणा-संकारक $[!B]$.
\end{enumerate}
\end{defn}

सार्वजनिक घोषणा-संकारक बॉक्स-संकारकासारखे कार्य
करतो. त्यानुसार भाषेची आगमनाधारित व्याख्या पुढीलप्रमाणे:

\begin{defn}
  ज्ञानविषयक भाषेतील \emph{!!^{formula}s} ची पुढीलप्रमाणे
  आगमनाने व्याख्या करतो:
  \begin{enumerate}
  \tagitem{prvFalse}{$\lfalse$ हे अण्विक !!{formula} आहे.}{}

  \tagitem{prvTrue}{$\ltrue$ हे अण्विक !!{formula} आहे.}{}

  \item प्रत्येक विधानीय चल~$\Obj p_i$ हे (अण्विक)
  !!{formula} आहे.

  \tagitem{prvNot}{$!A$ हे !!a{formula} असेल, तर $\lnot !A$
  हेही !!a{formula} आहे.}{}

  \tagitem{prvAnd}{$!A$ आणि $!B$ ही !!{formula}s असतील, तर $(!A \land
  !B)$ हे !!a{formula} आहे.}{}

  \tagitem{prvOr}{$!A$ आणि $!B$ ही !!{formula}s असतील, तर $(!A \lor !B)$
  हे !!a{formula} आहे.}{}

  \tagitem{prvIf}{$!A$ आणि $!B$ ही !!{formula}s असतील, तर $(!A \lif !B)$
  हे !!a{formula} आहे.}{}

  \tagitem{prvIff}{$!A$ आणि $!B$ ही !!{formula}s असतील, तर $(!A \liff !B)$
  हे !!a{formula} आहे.}{}

  \item $!A$ हे !!a{formula} असेल आणि $a \in G$ असेल,
  तर $\Knows_a !A$ हे !!a{formula} आहे.

  \item $!A$ आणि $!B$ ही !!{formula}s असतील,
  तर $[!A] !B$ हे !!a{formula} आहे.

  \tagitem{limitClause}{याखेरीज दुसरे काहीही !!a{formula} नाही.}{}
\end{enumerate}
\end{defn}

!!{formula} $[!A] !B$ चा अभिप्रेत अर्थ ``$!A$ सत्य
असल्याची सार्वजनिक घोषणा झाल्यानंतर $!B$~सत्य असते''
असा आहे. सार्वजनिक घोषणांच्या संदर्भात सामायिक
ज्ञानाविषयी बोलणेही कधी उपयोगी ठरते; म्हणून भाषेत
सामायिक-ज्ञान संकारक~$\CKnows_G
!A$ देखील समाविष्ट करता येतो.

\end{document}
